% 无穷矩阵 a_{i,j}=1{j=i}-1{j=i+1} 左上角 5 行 6 列；行和均为 0，第 1 列和为 1，其余列和为 0（第 6 列只显示前 5 行中的 -1，第 6 行的 +1 在图外）。
\begin{tikzpicture}[
  cell/.style={draw=FigureGrid,line width=.5pt,minimum size=7mm,inner sep=0pt},
  pcell/.style={cell,fill=FigureBlueFill,draw=FigureBlue},
  neg/.style={cell,fill=FigureRedFill,draw=FigureRed},
  note/.style={font=\FigureNoteFont\fontsize{7.5}{9}\selectfont,text=FigureMuted}]
\foreach \i in {1,...,5}{
  \foreach \j in {1,...,6}{
    \pgfmathtruncatemacro{\d}{\j-\i}
    \ifnum\d=0 \node[pcell] at (\j*.7,-\i*.7) {$+1$};
    \else\ifnum\d=1 \node[neg] at (\j*.7,-\i*.7) {$-1$};
    \else \node[cell] at (\j*.7,-\i*.7) {$0$};\fi\fi}
  \node at (6*.7+.95,-\i*.7) {$0$};
  \node at (.15,-\i*.7) {$i{=}\i$};}
\foreach \j in {1,...,6}{\node at (\j*.7,-.25) {$j{=}\j$};}
\node[note] at (6*.7+.95,-.25) {行和};
\foreach \j/\s in {1/1,2/0,3/0,4/0,5/0,6/0}{\node at (\j*.7,-6*.7+.05) {$\s$};}
\node[note,anchor=east] at (.45,-6*.7+.05) {列和};
\node[note] at (3.5*.7+.35,-6*.7-.5) {矩阵继续向右下方延伸};
\end{tikzpicture}
